%%%%%%%%%%%%%%%%%%%%%%%%%%%%
\label{sec:pseudocode}
%%%%%%%%%%%%%%%%%%%%%%%%%%%%
% For clarity, we introduce a few notations for the sets of polynomials, arising at the various stages of a $\deg_i$ iteration.
For $1 \le i \le d-1$, let $\mathcal{D}_i$ denote the set of polynomials obtained after the diagonalization at iteration-$i$ and $\mathcal{D}_d$ denotes the first diagonalized set.
For $1 \le i \le d-1$, let $\mathcal{P}_i$ denote the set of $S$-polynomials generated from $\mathcal{D}_{i+1}$, where $\mathcal{P}_d$ denotes the initial sample set. For each iteration-$i$, these sets have the same cardinality:
$|\mathcal{P}_i| = |\mathcal{D}_i| = \binom{n+i-1}{i}$.


%%%%%%%%%%%%%%%%%%%%%%%%%%%%%
%%%%%%%%%%%%%%%%%%%%%%%%%%%%%


\Cref{alg:step1} generates a set of $S$-polynomials $\mathcal{P}_i$ from the set of the diagonalized polynomials $\mathcal{D}_{i+1}$. This process selects a number of
admissible pairs (\Cref{def:admissible-pair}) $(f,g)$ to compute $S_{fg}$ (having degree $i+1$), and removes all monomials of degree $i+1$ by reducing w.r.t.\ the polynomial set $\mathcal{D}_{i+1}$.
\algdiag\ transforms a set of polynomials $\mathcal{P}_i$ into diagonal form (Smith Normal Form) $\mathcal{D}_i$.
% The main algorithm, \Cref{alg:alg}, iteratively applies $S$-polynomial generation and diagonalization steps to reduce the initial set of polynomials $\mathcal{P}_d$ to a set of $n$ linear polynomials of the form $\{x_j - s_j\} _{j=1}^n$.


%%%%%%%%%%%%%%%%%%%%%%%%%%%%%
%%%%%%%%%%%%%%%%%%%%%%%%%%%%%

\subsection{S-polynomial generation and reduction}\label{sec:onspoly}
\begin{definition}[Admissible Pair]\label{def:admissible-pair}
    Given two tuples of integers \(\alpha, \beta \in \mathbb{Z}_{\ge 0}^n\) such that \(|\alpha|=|\beta|=t\), the pair \((\alpha, \beta)\) is called an \emph{admissible pair} if and only if there exist two distinct standard basis vectors \(e_j, e_k \in \mathbb{Z}_{\ge 0}^n\) with \(j \neq k\) such that \(\alpha_j > 0\) and \(\beta = \alpha - e_j + e_k\). Equivalently, we say that two polynomials \(f, g\) with \(\lm(f) = \vars{x}^\alpha\) and \(\lm(g) = \vars{x}^\beta\) form an \emph{admissible pair} of polynomials if and only if the pair \((\alpha, \beta)\) is admissible. In other words, if \(|\alpha|=|\beta|=t\), then \(\deg(\gcd(\lm(f),\lm(g))) = t-1\).
\end{definition}
% The $S$-polynomial generation (in \Cref{alg:step1}) selects admissible pairs (\cref{def:admissible-pair}) of polynomials from the set $\mathcal{D}_{i+1}$ and computes their corresponding $S$-polynomials. Then the $S$-polynomials are reduced by removing all monomials of degree $i+1$ to produce $\mathcal{P}_i$.

% \begin{algorithm}[t]
% \caption{(\texttt{Spoly}) --- $S$-polynomials generation}\label{alg:step1}
% \begin{algorithmic}[1]
% \Statex{\hspace*{-\algorithmicindent} \textbf{Input:} $\mathcal{D}_{i+1}$: diagonal polynomial set, $n \ge 1$: dimension, $1\le i\le d-1$: current degree}
% \Statex{\hspace*{-\algorithmicindent} \textbf{Output:} $\mathcal{P}_i$: set of reduced $S$-polynomials of degree $i$}
% \Statex{}
% \State{$\mathcal{P}_i \gets \emptyset$}
% \State\texttt{PAIRS} $\gets$ admissible pairs from $\mathcal{D}_{i+1}$ \Comment{see~\Cref{def:admissible-pair} and \Cref{lemma:spoly_number}}
% \For{each $(f,g) \in \texttt{PAIRS}$}
%     \State{$S_{fg} \gets S_{fg}$}
%     \State{Remove all monomials of degree $i+1$ from $S_{fg}$}
%     \State{$\mathcal{P}_i \gets \mathcal{P}_i \cup \{S_{fg}\}$}
% \EndFor{}
% \State{\Return{$\mathcal{P}_i$}}
% \end{algorithmic}
% \end{algorithm}

\begin{algorithm}[b]
\caption{(\texttt{NextAdmissiblePair}) --- Next admissible pair generation}\label{alg:next-admissible-pair}
\begin{algorithmic}[1]
\Statex{\hspace*{-\algorithmicindent} \textbf{Input:} $\ell_j$: $n$-tuple, last component of previous pair, $n \ge 1$: dimension, $2 \le i\le d$: current degree}
\Statex{\hspace*{-\algorithmicindent} \textbf{Output:} $(\ell_{j+1}, \ell_{j+2})$: next admissible pair}
\Statex{\hspace*{-\algorithmicindent} \textbf{Uses:} \textsc{Successor}: get next combination in Gray code order}\Comment{See Appendix~\ref{sec:successor}}
% ~\cite{DBLP:journals/cj/Takaoka99,DBLP:journals/corr/Takaoka15a}
% \If{$\ell_j == \varnothing$}
%     \State{$\ell_{j+1} \gets$ (0, 0, \ldots, 0, $i$)} \Comment{First combination in Gray code order~\cite{DBLP:journals/ejc/RuskeyS96}}
% \Else
    \State{$\ell_{j+1} \gets$ \textsc{Successor}($\ell_j$, $n$, $i$)} 
% \EndIf
\State{$\ell_{j+2} \gets$ \textsc{Successor}($\ell_{j+1}$, $n$, $i$)} 
% \State{\texttt{PAIR} $\gets$ $(\ell_{j+1}, \ell_{j+2})$}
% \State{\Return{\texttt{PAIR}}}
\State{\Return{$(\ell_{j+1}, \ell_{j+2})$}}
\end{algorithmic}
\end{algorithm}

\begin{algorithm}[t]
\caption{(\texttt{Spoly}) --- $S$-polynomials generation}\label{alg:step1}
\begin{algorithmic}[1]
\Statex{\hspace*{-\algorithmicindent} \textbf{Input:} $\mathcal{D}_{i+1}$: diagonal polynomial set, $n \ge 1$: dimension, $1\le i\le d-1$: current degree}
\Statex{\hspace*{-\algorithmicindent} \textbf{Output:} $\mathcal{P}_i$: set of reduced $S$-polynomials of degree $i$}

\State{$\mathcal{P}_i \gets \emptyset$}\Comment{Initialize the set of $S$-polynomials}
\State{PrevCombination $\gets \varnothing$}
% \State\texttt{PAIRS} $\gets$ admissible pairs from $\mathcal{D}_{i+1}$ \Comment{see~\Cref{def:admissible-pair} and \Cref{lemma:spoly_number}}
\For{$k\gets1$ to $\binom{n+i-1}{i}$ by $1$}
    \State{$(\ell_{2k-1}, \ell_{2k}) \gets$ \texttt{NextAdmissiblePair}(PrevCombination, $n$, $i+1$)} 
    \State{PrevCombination $\gets \ell_{2k}$}
    % \State{Select $(\ell_{2k-1}, \ell_{2k})$} \Comment{see~\Cref{def:admissible-pair} and \Cref{lemma:spoly_number}}
    \State{Select $(f,g) \in \mathcal{D}_{i+1}$ such that \(\LM(f) = \vars{x}^{{\ell}_{2k-1}}\) and \(\LM(g) = \vars{x}^{\ell_{2k}}\)} 
    \State{Compute $S_{fg}$} \Comment{See \Cref{eq:spoly}}
    \State{Reduce $S_{fg}$ with the polynomials in $\mathcal{D}_{i+1}$} 
    \State{$\mathcal{P}_i \gets \mathcal{P}_i \cup \{S_{fg}\}$}
\EndFor{}
\State{\Return{$\mathcal{P}_i$}}
\end{algorithmic}
\end{algorithm}

% We now specify the criterion used to select admissible pairs of polynomials for $S$-polynomial generation and explain the motivation behind this choice.
%\paragraph*{}\label{paragraph:criterion}
Consider the iteration-$i$. Note that given a monomial $m$ of degree \(i+1\), due to the diagonal form of $\mathcal{D}_{i+1}$ there exist a unique polynomial $f \in \mathcal{D}_{i+1}$ such that \(\lm(f) = m\). Hence, given two monomials $m_1, m_2$ that form an admissible pair, the corresponding admissible pair of polynomials can be uniquely determined. Let \(f,g\) be those two polynomials, the resulting $S$-polynomial \(S_{fg}\) admits the simplified form given in \Cref{eq:spoly}.

Since they form an admissible pair, the least common multiple of the leading terms has degree $i+2$.
Consequently, recalling that $\deg(\lm(f))=\deg(\lm(g))=i+1$, both $\frac{\vars{x}^\gamma}{\lm(f)}$ and $\frac{\vars{x}^\gamma}{\lm(g)}$ (in \Cref{eq:spoly-expanded}) are monomials of degree one. Using the notation of Section~\ref{sec:2background}, we decompose
\[
    f = f^{(i+1)} + f^{(i)} + \cdots + f^{(1)} + f^{(0)}.
\] 
Each polynomial in $\mathcal{D}_{i+1}$ contains a unique monomial of degree $i+1$, namely $\lm(f)$.
Multiplication by a variable $x_k$ therefore yields
\[
    x_k \cdot f = f^{(i+2)} + f^{(i+1)} + \cdots + f^{(1)},
\]
where $f^{(i+2)}$ consists of a single monomial, namely $x_k \cdot \lm(f)$.
As a result, when forming the $S$-polynomial of two such polynomials, the degree-$(i+2)$ terms cancel, and we obtain
\[
    S_{fg} = {S_{fg}}^{(i+1)} + {S_{fg}}^{(i)} + \cdots + {S_{fg}}^{(2)} + {S_{fg}}^{(1)}.
\]
Thus, $S_{fg}$ contains only monomials of degree at most $i+1$.
Since the leading terms of the polynomials in $\mathcal{D}_{i+1}$ exhaust all monomials of degree $i+1$, these terms can be eliminated by reduction with respect to $\mathcal{D}_{i+1}$, yielding a polynomial of degree at most $i$.

Consequently, this reduction involves only linear operations and no divisions.
This property is guaranteed by the chosen criterion and ensures that the reduced $S$-polynomial has the form
\[
    S_{fg} = {S_{fg}}^{(i)} + {S_{fg}}^{(i-1)} + \cdots + {S_{fg}}^{(1)} + {S_{fg}}^{(0)}.
\]


%%%%%%%%%%%%%%%%%%%%%%%%%%%%%
%%%%%%%%%%%%%%%%%%%%%%%%%%%%%

\noindent
\textbf{Existence and construction of admissible pairs.} 
\Cref{lemma:spoly_number} guarantees the existence of a sufficient number of admissible pairs to generate all required $S$-polynomials, while preventing the formation of redundant pairs.
% or linearly dependent $S$-polynomials.

\begin{lemma}[Feasibility condition]\label{lemma:spoly_number}
% Given a degree $2\le i< n$, and \(\mathcal{D}_i\) (a set of $\binom{n+i-1}{i}$ diagonals polynomials of degree $i$), it is always possible to group them in at least $\binom{n+i-2}{i-1}$ distinct admissible pairs.

Let $1 \le i < n-1$, and let $\mathcal{D}_{i+1}$ be the set of diagonal polynomials of degree $i+1$, with $|\mathcal{D}_{i+1}| = \binom{n+i}{i+1}$.
Then there exists a collection of (at least) $\binom{n+i-1}{i}$ pairwise disjoint admissible pairs of elements of $\mathcal{D}_{i+1}$.

%Following the described criterion, \Cref{alg:pairs-construction} generates at least \(\binom{n+d-2}{d-1}\) \(S\)-polynomials, starting from $\mathcal{D}_{i+1}$ at $\deg_i$ iteration.
% Given a degree \(d\), a number of variables \(n\), and a polynomial set of size \(\binom{n+d-1}{d}\) whose leading monomials exhaust all monomials of degree \(d\), \Cref{alg:pairs-construction} generates at least \(\binom{n+d-2}{d-1}\) \(S\)-polynomials if and only if \(d \le n-1\).
\end{lemma}

\begin{proof}
In order to prove this result, we first establish \emph{(i)} a necessary condition for the existence of a sufficient number of admissible pairs, and then show \emph{(ii)} that this condition is also sufficient by providing an explicit construction of such pairs.


\begin{enumerate}[\itshape i)]
    \item To generate at least \(\binom{n+i-1}{i}\) \emph{pairwise disjoint} pairs, at least \(2\binom{n+i-1}{i}\) are required. 
    Hence, a necessary condition is
    \[
    \begin{split}
       \binom{n + i}{i+1} &\ge 2\binom{n+i-1}{i} \implies \frac{(n+i)!}{(n-1)!(i+1)!} \ge 2\frac{(n+i-1)!}{(n-1)!(i)!}\\
       \implies \frac{n+i}{i+1} &\ge 2 \implies n+i \ge 2(i+1) \implies  n \ge i+2 \implies n > i+1
    \end{split}
    \] where \(i+1\), in the analyzed problem, can be at maximum \(d\). As a result, if \(n > d\), we can generate enough distinct pairs. 

    % \item Consider a degree $i$ such that $2\le i < n$, and the corresponding set of diagonal polynomials $D_{i}$. We construct the set \(\mathcal{A} := \{ \alpha \colon \exists f \in \mathcal{D}_i \ \left(\LM(f) = \vars{x}^\alpha \right)\}\). Given \(\alpha, \beta \in \mathcal{A}\), we say that \(\alpha\) and \(\beta\) are an \emph{admissible pair} if and only if there exist two standard basis vectors \(e_j, e_k\) where \(j \neq k\) such that \(\alpha_j > 0\) and \(\beta = \alpha - e_j + e_k\). This configuration was already considered by~\cite{DBLP:journals/ejc/RuskeyS96}. More specifically, \(\mathcal{A}\) can be seen as the set of all combinations of \(n\) elements chosen \(i\) at a time. In other words the elements of \set{A} can be regarded as placements of \(i\) identical balls into \(n\) distinct boxes where each box can hold at most \(i\) balls. As a result, the elements of \set{A} can be listed in such a way that consecutive elements form admissible pairs. Such listing is known as \emph{Gray Code for combinations of a multiset}~\cite{DBLP:journals/ejc/RuskeyS96} and it can be interpreted as a Hamiltonian path in the \emph{graph} where the vertices are the elements of \set{A} and two vertices are connected if and only if they form an admissible pair. By definition, a Hamiltonian path covers all the vertices of the graph, hence we can generate at least \(\left\lfloor\frac{\binom{n+i-1}{i}}{2}\right\rfloor\) distinct admissible pairs. Let \(\ell_j\) for \(1 \le j \le \binom{n+i-1}{i}\) be the elements in the constructed Hamiltonian path, we can construct \((\ell_{2k-1}, \ell_{2k})\) for \(1 \le k \le \left\lfloor\frac{\binom{n+i-1}{i}}{2}\right\rfloor\). This way, as shown in \textbf{\textsf{(i)}}, when \(n > d\) we can generate at least \(\binom{n+i-2}{i-1}\) distinct admissible pairs.

    \item Let $1 \le i < n-1$ and let $\mathcal{D}_{i+1}$ be the set of diagonal polynomials of degree $i+1$.
    We define $\mathcal{A} := \{ \alpha \in \mathbb{Z}_{\ge 0}^n \colon \exists f \in \mathcal{D}_{i+1} \text{ with } \lm(f) = \vars{x}^\alpha \}$ i.e., the set of all multi-indices corresponding to the leading monomials of polynomials in $\mathcal{D}_{i+1}$.
    % For $\alpha, \beta \in \mathcal{A}$, we say that $(\alpha,\beta)$ is an \emph{admissible pair} if and only if there exist two distinct standard basis vectors $e_j, e_k$ such that $\alpha_j > 0$ and $\beta = \alpha - e_j + e_k$. 
    Following~\cite{DBLP:journals/ejc/RuskeyS96}, we observe that $\mathcal{A}$ can be viewed as the set of all combinations of $n$ elements chosen $i+1$ at a time, or as placements of $i+1$ identical balls into $n$ distinct boxes, with each box containing at most $i+1$ balls. In this perspective, the elements of $\mathcal{A}$ can be ordered so that consecutive elements form admissible pairs. Such an ordering is known as a \emph{Gray code for combinations of a multiset}~\cite{DBLP:journals/ejc/RuskeyS96}.  
    This Gray code can be interpreted as a Hamiltonian path in the graph whose vertices are the elements of $\mathcal{A}$, with edges connecting pairs of vertices if and only if they form an admissible pair. By definition, a Hamiltonian path covers all vertices of the graph, which allows us to generate at most $\left\lfloor {\binom{n+i}{i+1}}/{2} \right\rfloor$ distinct admissible pairs. Denoting the elements along the path by $\ell_j$, $1 \le j \le \binom{n+i}{i+1}$, we may form the pairs $(\ell_{2k-1}, \ell_{2k})$ for $1 \le k \le \left\lfloor {\binom{n+i}{i+1}}/{2} \right\rfloor$.
    This ensures that each polynomial appears in at most one pair and that all pairs are admissible.
    Hence, as shown in \textcolor{gray}{\textbf{\textsf{(i)}}}, when $n>i+1$ we can generate (at least) $\binom{n+i-1}{i}$ pairwise distinct \emph{admissible} pairs.
\end{enumerate}
\end{proof} 
% As shown in~\cite{DBLP:journals/cj/Takaoka99,DBLP:journals/corr/Takaoka15a}, the Gray code for combinations of a multiset~\cite{DBLP:journals/ejc/RuskeyS96} can be generated efficiently, either via a recursive procedure or its iterative variant. Moreover, these results provide a method to determine the successor of any given combination in constant time. Consequently, the admissible pairs can be generated alongside the corresponding $S$-polynomial computations, without increasing the overall computational complexity of \Cref{alg:step1}.

% Gray codes for combinations of a multiset can be generated efficiently, either via recursive or iterative procedures, as shown in~\cite{DBLP:journals/cj/Takaoka99,DBLP:journals/corr/Takaoka15a,DBLP:journals/ejc/RuskeyS96}. In particular, these constructions allow the successor of a given combination to be computed in constant time. In our algorithm, the \texttt{Successor} functionality is used to generate admissible pairs on the fly during the $S$-polynomial computations, without increasing the overall computational complexity of \Cref{alg:step1}. Details on the adapted successor procedure used in this work are provided in \Cref{alg:successor} in Appendix \ref{sec:successor}.


\emph{Gray codes for combinations of a multiset}~\cite{DBLP:journals/ejc/RuskeyS96} admit a constant-time successor computation algorithm~\cite{DBLP:journals/cj/Takaoka99,DBLP:journals/corr/Takaoka15a}. This successor procedure can be used to efficiently determine the Hamiltonian path of the adjacency graph described in the proof of \Cref{lemma:spoly_number}. Consequently, admissible pairs can be generated incrementally and on-demand during the $S$-polynomial generation phase, without increasing the asymptotic complexity of~\Cref{alg:step1}.
For the sake of completeness, a pseudocode description of the \textsc{Successor} procedure is given as \Cref{alg:successor} in Appendix~\ref{sec:successor}.

%%%%%%%%%%%%%%%%%%%%%%%%%%%%%
%%%%%%%%%%%%%%%%%%%%%%%%%%%%%

% \noindent
% \textbf{\textsf{Diagonalization and reduction.}} 




% \begin{algorithm}[t!]
% \caption{(\algdiag\) --- Diagonalization: conversion to Smith Normal Form}\label{alg:step2b}
% \begin{algorithmic}[1]
% \Statex{\hspace*{-\algorithmicindent} \textbf{Input:} $\mathcal{P}_i$: set of polynomials, $n \geq 1$: secret's dimension, $1\leq i\leq d$: current degree}
% \Statex{\hspace*{-\algorithmicindent} \textbf{Output:} $\mathcal{D}_i$: set of polynomials in diagonal form}

% \State{$\mathcal{D}_i \gets \texttt{HermiteForm}(\mathcal{P}_i)$}

% \State{\Return{$\mathcal{D}_i$}}
% \end{algorithmic}
% \end{algorithm}


%%%%%%%%%%%%%%%%%%%%%%%%%%%%%
%%%%%%%%%%%%%%%%%%%%%%%%%%%%%

\subsection{Computational Complexity}
\begin{theorem}\label{thm:iterationCost}
    In \Cref{alg:alg} the complexities of $\algspoly$, $\algdiag$ are given as follows:
    \begin{enumerate}[\itshape i)]
        \item $\mathscr{C}_{\algspoly} = \mathcal{O}\left( n^{3(d-1)} \right)$ over $d-1$ iterations,
        \item $\mathscr{C}_{\algdiag} = \bigO\left(n^{d\omega}\right)$, over $d$ iterations. 
        %\item $C_{\texttt{DiagB}} = \mathcal{O}\!\left(\dfrac{dn^2}{2{(n+d)}^3}\dbinom{n+d}{d}^3\right)$, over $d$ iterations.
    \end{enumerate}
    Hence, the overall complexity of \Cref{alg:alg} is $\mathcal{O}\left(n^{\max(3(d-1), d\omega)}\right)$ scalar operations.
    % The overall complexity of \Cref{alg:alg} is given by the sum of the complexities (over $1\le i\le d-1$ for \texttt{Spoly} and $1\le i\le d$ for \texttt{DiagA} and \texttt{DiagB}) of:
    % \begin{itemize}
    %     \item $S$-polynomial generation and reduction:
    %         $C_{\texttt{Spoly}}(i) = 2\cdot \binom{n+i-1}{i}^2\cdot \!\left[ \binom{n+i}{i} + 1 \right]$,
    %     \item conversion to Smith Normal Form: $C_{\texttt{DiagA}}(i) =\sum_{j=1}^{\binom{n+i-1}{i}}\!\left[\!\left[ \binom{n+i}{i} - j + 1 \right]\cdot \!\left[ \binom{n+i-1}{i} - j + 1 \right]\right]$,
    %     \item conversion to diagonal form: $C_{\texttt{DiagB}}(i) =\sum_{j=1}^{\binom{n+i-1}{i}} \!\left[\!\left[ \binom{n+i-1}{i-1} +1 \right] \cdot \!\left[ \binom{n+i-1}{i} - j \right]\right]$.
    % \end{itemize}
\end{theorem}

%The proofs of the \cref{thm:step1,thm:step2a,thm:step2b} are given in \Cref{sec:cc}.
\begin{proof}
    % Follows directly from the complexity analyses of the individual steps, provided in \Cref{thm:step1,thm:step2a,thm:step2b} in \Cref{sec:cc}.
    The complete proof is given in \Cref{sec:complexity}.    
\end{proof}

% As stated in \Cref{thm:result}, the complexity of \Cref{alg:alg} is $\mathcal{O}\!\left(\frac{1}{3}\binom{n+d}{d}^3\right)$ field operations.

\begin{proof}[Proof of \Cref{thm:result}]
    The complexity $\mathscr{C}_{tot}$ of \Cref{alg:alg} follows immediately by adding the complexities of the two main procedures, as given in \Cref{thm:iterationCost}, i.e.
    \begin{align*}
        \mathscr{C}_{tot} &:= \mathscr{C}_{\algspoly} + \mathscr{C}_{\algdiag} = \mathcal{O}\left( n^{3(d-1)} \right) + \mathcal{O}\left(n^{d\omega}\right) = \mathcal{O}\left(n^{\max(3(d-1), d\omega)}\right).
    \end{align*}
    The correctness of the algorithm is proven in \Cref{sec:correct}.
\end{proof}


%%%%%%%%%%%%%%%%%%%%%%%%%%%%%
%%%%%%%%%%%%%%%%%%%%%%%%%%%%%
% \noindent
% \textbf{\belwe.} For this case, our generic analysis for degree $d$ can be applied directly with $d = 2$. Thus, we require $\mathcal{O}\left(\frac{n(n+1)}{2}\right)$ samples and the complexity is
% \[
%   \mathcal{O}\left(\binom{n+2}{2}^3 \right) = \mathcal{O}\left(\frac{1}{8}\cdot n^6 \right).
% \]

% \noindent
% \textbf{\bslwe.} For this case, we add the constraint set $\{ x_j^2 - x_j = 0: 1 \le j \le n\}$ to the polynomial system. Consequently, all monomials are of the form $x_1^{\alpha_1}\cdots x_n^{\alpha_n}$, where $\alpha_j \in \{0,1\}$ for ${1\le j\le n}$. The number of samples required in this case is $\binom{n}{d}$ and the complexity is $\mathcal{O}(\binom{n}{d}^3)$. A slighlty more detailed discussion is given in the Appendix~\ref{sec:binarysecret}.
% %$\mathcal{O}((^nC_k)^3)$

\subsection{Correctness}\label{sec:correct}

Throughout this section, we assume that the error terms satisfy the prescribed bound. Thus, the secret $\vars{s}$ is a common zero of all polynomials in the initial set of samples $\mathcal{P}_d$.

\begin{proposition}[Preservation of the LWE secret]\label{prop:preservation}
At each iteration-$i$ of \Cref{alg:alg}, with $1 \le i \le d-1$, the secret $\vars{s}$ is a common zero of all polynomials in $\mathcal{D}_i$.
\end{proposition}

\begin{proof}
\textbf{Base case.}  
$\mathcal{D}_d$ is obtained from $\mathcal{P}_d$ via diagonalization (\algdiag), which consists of linear combinations of polynomials in $\mathcal{P}_d$.  
Since $\vars{s}$ is a common zero of all polynomials in $\mathcal{P}_d$, it is also a zero of all polynomials in $\mathcal{D}_d$.

We proceed by induction on $i$, decreasing from $i=d-1$ to $i=1$.

\noindent
\textbf{Inductive step.}  
% Assume that $\vars{s}$ is a common zero of all polynomials in $\mathcal{D}_{i+1}$, for $1 \le i \le d-1$. Consider the iteration-$i$:
Assume that, at iteration-$i$, $\vars{s}$ is a common zero of all polynomials in $\mathcal{D}_{i+1}$.
The following steps are performed:
\begin{enumerate}
    \item \emph{$S$-polynomial generation.} Let $f,g\in\mathcal{D}_{i+1}$. By the inductive hypothesis, they vanish at $\vars{s}$. By definition (\Cref{eq:spoly-expanded}), also their $S$-polynomial vanishes at $\vars{s}$.

    \item \emph{$S$-polynomial reduction, and diagonalization.} We showed in \Cref{sec:onspoly} that reductions do not involve polynomial divisions.
    Therefore, these steps only consist of linear combinations of previously
    generated polynomials. Since all such polynomials vanish at
    $\vars{s}$, the resulting polynomials also vanish at $\vars{s}$.
\end{enumerate}
Therefore, $\vars{s}$ is a common zero of all polynomials in $\mathcal{D}_i$, for all $1 \le i \le d$.
\end{proof}

%%%%%%%%%%%%%%%%%%%%%%%%%%%%%
%%%%%%%%%%%%%%%%%%%%%%%%%%%%%


% \begin{corollary}[Recovery of the secret]\label{cor:secret_recovery}
% At the last iteration, the secret $\vars{s}$ is a common zero of the linear equations in $\mathcal{D}_1$.
% Each equation, due to diagonalization process, is of the form $ x_i - s_i$ for $i \in \{1,\dots,n\}$, and the secret is uniquely determined as $\vars{s} = (s_1, \dots, s_n)$.
% \end{corollary}

% \begin{proof}
% By \Cref{prop:preservation}, the property of vanishing at $\vars{s}$ is preserved throughout all iterations.
% At the last iteration, all polynomials have degree $1$, and in particular they are linear and monic.
% Each polynomial is therefore of the form $x_i - s_i$ for some $i$, and the secret vector $\vars{s}$ is uniquely recovered as $\vars{s} = (s_1, \dots, s_n)$.
% \end{proof}


By \Cref{prop:preservation}, $\vars{s}$ remains the zero of the polynomials through all iterations.
In particular, at the last iteration, all polynomials in $\mathcal{D}_1$ are linear and of degree $1$.
Each equation, due to diagonalization process, is of the form $a_i x_i -b_i$ for $i \in \{1,\dots,n\}$, and the secret can be determined as described in \Cref{subsec:highlevel}.

%%%%%%%%%%%%%%%%%%%%%%%%%%%%%
%%%%%%%%%%%%%%%%%%%%%%%%%%%%%


\begin{proposition}[Termination]\label{prop:termination}
\Cref{alg:alg} terminates after a finite number of iterations.
\end{proposition}

\begin{proof}
  At each iteration of the algorithm the total degree is decreased by $1$. Thus, the algorithm stops after $d-1$ iterations, yielding linear equations. 
% Consider the iteration-$i$ of \Cref{alg:alg}.  
% At this iteration, \texttt{Spoly} (\Cref{alg:step1}) generates $S$-polynomials of degree at most $i$ from the degree-$(i+1)$ polynomials in $\mathcal{D}_{i+1}$.  
% Consequently, the maximum total degree of the polynomials strictly decreases at each iteration.  
% Since the degree is a non-negative integer and the algorithm halts once degree one is reached, the process terminates after at most $d$ iterations.
\end{proof}


%%%%%%%%%%%%%%%%%%%%%%%%%%%%%
%%%%%%%%%%%%%%%%%%%%%%%%%%%%%


% \begin{remark}[On the solution set]
% \Cref{alg:alg} does not, in general, preserve the full solution set of the LWE polynomial system.  
% Its design ensures that the LWE secret $\vars{s}$ is preserved, which is the unique solution of interest for the \slwe\ problem.
% \end{remark}


